% =====================================================================
%  minggu-12.tex  --  Slide Pertemuan Minggu ke-12
%  Pengantar Kriptografi (MPFA014) -- Prodi Matematika FMIPA UBB
%
%  Topik silabus : Diffie--Hellman dan ElGamal
%  Sumber materi : buku_ajar/chapters/11 (Bab 11), 12 (Bab 12)
%  Sub-CPMK      : 1.4
%
%  Kompilasi dari dalam folder ppt/:
%      pdflatex -interaction=nonstopmode minggu-12.tex   (dua kali)
% =====================================================================

\documentclass[aspectratio=169,11pt]{beamer}
% =====================================================================
%  beamer-preamble.tex  --  prasarana dan gaya slide kuliah
%
%  Mata kuliah : Pengantar Kriptografi (MPFA014)
%  Program     : Program Studi Matematika, FMIPA, Universitas Bale Bandung
%
%  Berkas ini TIDAK memuat \documentclass. Setiap berkas minggu memuatnya
%  sendiri lalu menarik berkas ini:
%
%      \documentclass[aspectratio=169,11pt]{beamer}
%      % =====================================================================
%  beamer-preamble.tex  --  prasarana dan gaya slide kuliah
%
%  Mata kuliah : Pengantar Kriptografi (MPFA014)
%  Program     : Program Studi Matematika, FMIPA, Universitas Bale Bandung
%
%  Berkas ini TIDAK memuat \documentclass. Setiap berkas minggu memuatnya
%  sendiri lalu menarik berkas ini:
%
%      \documentclass[aspectratio=169,11pt]{beamer}
%      % =====================================================================
%  beamer-preamble.tex  --  prasarana dan gaya slide kuliah
%
%  Mata kuliah : Pengantar Kriptografi (MPFA014)
%  Program     : Program Studi Matematika, FMIPA, Universitas Bale Bandung
%
%  Berkas ini TIDAK memuat \documentclass. Setiap berkas minggu memuatnya
%  sendiri lalu menarik berkas ini:
%
%      \documentclass[aspectratio=169,11pt]{beamer}
%      \input{beamer-preamble}
%      \title{...} \subtitle{...} \setWeekInfo{Minggu n}{Topik}
%
%  Berkas minggu hanya berisi \setWeekInfo dan lingkungan frame.
%  Kompilasi dijalankan DARI DALAM folder ppt/ (lihat compile-ppt.bat),
%  karena gambar buku dirujuk dengan jalur relatif ../buku_ajar/.
% =====================================================================

% ---------------------------------------------------------------- bahasa
\usepackage[utf8]{inputenc}
\usepackage[T1]{fontenc}
\usepackage[indonesian]{babel}
\usepackage{lmodern}
\usepackage{microtype}

% ------------------------------------------------------------- matematika
\usepackage{amsmath,amssymb,mathtools}

% ----------------------------------------------------------------- tabel
\usepackage{booktabs}
\usepackage{array}
\usepackage{tabularx}

% --------------------------------------------------------------- grafik
\usepackage{tikz}
\usetikzlibrary{arrows.meta,positioning,shapes.geometric,calc,fit,backgrounds}

% ---------------------------------------------------------- kode program
\usepackage{listings}
\lstdefinestyle{slidecode}{
    basicstyle=\ttfamily\scriptsize,
    keywordstyle=\color{ubbblue}\bfseries,
    commentstyle=\color{ubbgreen}\itshape,
    stringstyle=\color{ubbred},
    numbers=left,
    numberstyle=\tiny\color{gray},
    numbersep=6pt,
    frame=single,
    rulecolor=\color{ubbblue!40},
    backgroundcolor=\color{ubbblue!3},
    breaklines=true,
    showstringspaces=false,
    tabsize=4,
    columns=fullflexible,
    keepspaces=true,
}
\lstset{style=slidecode}

% =====================================================================
%  WARNA INSTITUSI
% =====================================================================
\definecolor{ubbblue}{RGB}{12, 62, 120}      % warna utama
\definecolor{ubbteal}{RGB}{0, 118, 110}      % aksen kedua
\definecolor{ubbgreen}{RGB}{22, 122, 74}     % untuk blok contoh
\definecolor{ubbred}{RGB}{168, 42, 42}       % untuk blok peringatan
\definecolor{ubbgray}{RGB}{70, 76, 84}       % teks sekunder

% =====================================================================
%  IDENTITAS MATA KULIAH
% =====================================================================
\newcommand{\CourseName}{Pengantar Kriptografi}
\newcommand{\CourseCode}{MPFA014}
\newcommand{\CourseSemester}{Semester Ganjil}
\newcommand{\CourseProgram}{Program Studi Matematika}
\newcommand{\CourseFaculty}{Fakultas Matematika dan Ilmu Pengetahuan Alam}
\newcommand{\CourseInstitution}{Universitas Bale Bandung}
\newcommand{\CourseLecturer}{Cecep Suwanda, S.Si., M.Kom.}

% Label minggu, disetel tiap berkas lewat \setWeekInfo
\newcommand{\WeekLabel}{}
\newcommand{\WeekTopic}{}
\newcommand{\setWeekInfo}[2]{%
    \renewcommand{\WeekLabel}{#1}%
    \renewcommand{\WeekTopic}{#2}%
}

% =====================================================================
%  TEMA: header, footer, warna blok
% =====================================================================
\usetheme{default}
\setbeamertemplate{navigation symbols}{}
\setbeamertemplate{blocks}[rounded][shadow=false]

\setbeamercolor{palette primary}{bg=ubbblue, fg=white}
\setbeamercolor{palette secondary}{bg=ubbgray, fg=white}
\setbeamercolor{palette tertiary}{bg=ubbteal, fg=white}

\setbeamercolor{structure}{fg=ubbblue}
\setbeamercolor{normal text}{fg=black!92}
\setbeamercolor{alerted text}{fg=ubbred}

\setbeamercolor{block title}{bg=ubbblue, fg=white}
\setbeamercolor{block body}{bg=ubbblue!6, fg=black!92}
\setbeamercolor{block title alerted}{bg=ubbred, fg=white}
\setbeamercolor{block body alerted}{bg=ubbred!6, fg=black!92}
\setbeamercolor{block title example}{bg=ubbgreen, fg=white}
\setbeamercolor{block body example}{bg=ubbgreen!6, fg=black!92}

\setbeamercolor{itemize item}{fg=ubbblue}
\setbeamercolor{itemize subitem}{fg=ubbteal}
\setbeamercolor{enumerate item}{fg=ubbblue}

\setbeamerfont{headline}{size=\scriptsize}
\setbeamerfont{footline}{size=\scriptsize}
\setbeamerfont{block title}{size=\normalsize}

% Badan blok rata kiri (bukan rata kanan-kiri): teks slide lebih mudah dibaca
% dan baris terakhir yang pendek tidak lagi memicu "Underfull \hbox".
\addtobeamertemplate{block body begin}{\raggedright}{}
\addtobeamertemplate{block alerted body begin}{\raggedright}{}
\addtobeamertemplate{block example body begin}{\raggedright}{}

% Header: nama mata kuliah dan kode MK
\setbeamertemplate{headline}{%
    \begin{beamercolorbox}[wd=\paperwidth,ht=2.6ex,dp=1.2ex,leftskip=1em,rightskip=1em]{palette primary}%
        \usebeamerfont{headline}\CourseName\ \textperiodcentered\ \CourseCode%
        \hfill \CourseProgram%
    \end{beamercolorbox}%
}

% Footer: minggu dan nomor halaman
\setbeamertemplate{footline}{%
    \begin{beamercolorbox}[wd=\paperwidth,ht=2.4ex,dp=1.2ex,leftskip=1em,rightskip=1em]{palette secondary}%
        \usebeamerfont{footline}\WeekLabel\ \textperiodcentered\ \WeekTopic%
        \hfill \insertframenumber\,/\,\inserttotalframenumber%
    \end{beamercolorbox}%
}

% Halaman judul: judul, subjudul, lalu blok identitas mata kuliah
\setbeamertemplate{title page}{%
    \begin{centering}
        \vskip1.0em
        {\usebeamerfont{title}\usebeamercolor[fg]{title}\inserttitle\par}
        \vskip0.3em
        {\usebeamerfont{subtitle}\usebeamercolor[fg]{subtitle}\insertsubtitle\par}
        \vskip1.2em
        \textcolor{ubbblue}{\rule{0.72\textwidth}{0.8pt}}
        \vskip1.2em
        {\small
            \textbf{\CourseName}\ (\CourseCode)\par
            \vskip0.4em
            \CourseProgram\par
            \CourseFaculty\par
            \CourseInstitution\par
            \vskip0.6em
            \CourseSemester\ \textperiodcentered\ \WeekLabel\par
            \vskip0.8em
            Dosen Pengampu: \CourseLecturer\par
        }
        \vskip1.0em
        \textcolor{ubbblue}{\rule{0.72\textwidth}{0.4pt}}
    \end{centering}
}

% =====================================================================
%  PERINTAH BANTU
% =====================================================================
\newcommand{\R}{\mathbb{R}}
\newcommand{\N}{\mathbb{N}}
\newcommand{\Z}{\mathbb{Z}}
\newcommand{\code}[1]{\texttt{#1}}

% Rujukan gambar dari buku: \bookfig[lebar]{NN}{namaberkas.png}
% Jalur relatif terhadap folder ppt/ (dikompilasi dari sana).
\newcommand{\bookfig}[3][0.80\linewidth]{%
    \includegraphics[width=#1]{../buku_ajar/chapters/#2/figures/#3}%
}

% Keterangan sumber gambar, sepadan dengan \sumber di buku
\newcommand{\figsumber}[1]{%
    \par\vspace{1pt}{\tiny\itshape Sumber: #1\par}%
}

% Penanda blok materi yang sering dipakai di slide
\newcommand{\slidejudul}[1]{{\large\bfseries\textcolor{ubbblue}{#1}}\par\vspace{2pt}}
\newcommand{\kunci}[1]{\textcolor{ubbred}{\textbf{#1}}}

%      \title{...} \subtitle{...} \setWeekInfo{Minggu n}{Topik}
%
%  Berkas minggu hanya berisi \setWeekInfo dan lingkungan frame.
%  Kompilasi dijalankan DARI DALAM folder ppt/ (lihat compile-ppt.bat),
%  karena gambar buku dirujuk dengan jalur relatif ../buku_ajar/.
% =====================================================================

% ---------------------------------------------------------------- bahasa
\usepackage[utf8]{inputenc}
\usepackage[T1]{fontenc}
\usepackage[indonesian]{babel}
\usepackage{lmodern}
\usepackage{microtype}

% ------------------------------------------------------------- matematika
\usepackage{amsmath,amssymb,mathtools}

% ----------------------------------------------------------------- tabel
\usepackage{booktabs}
\usepackage{array}
\usepackage{tabularx}

% --------------------------------------------------------------- grafik
\usepackage{tikz}
\usetikzlibrary{arrows.meta,positioning,shapes.geometric,calc,fit,backgrounds}

% ---------------------------------------------------------- kode program
\usepackage{listings}
\lstdefinestyle{slidecode}{
    basicstyle=\ttfamily\scriptsize,
    keywordstyle=\color{ubbblue}\bfseries,
    commentstyle=\color{ubbgreen}\itshape,
    stringstyle=\color{ubbred},
    numbers=left,
    numberstyle=\tiny\color{gray},
    numbersep=6pt,
    frame=single,
    rulecolor=\color{ubbblue!40},
    backgroundcolor=\color{ubbblue!3},
    breaklines=true,
    showstringspaces=false,
    tabsize=4,
    columns=fullflexible,
    keepspaces=true,
}
\lstset{style=slidecode}

% =====================================================================
%  WARNA INSTITUSI
% =====================================================================
\definecolor{ubbblue}{RGB}{12, 62, 120}      % warna utama
\definecolor{ubbteal}{RGB}{0, 118, 110}      % aksen kedua
\definecolor{ubbgreen}{RGB}{22, 122, 74}     % untuk blok contoh
\definecolor{ubbred}{RGB}{168, 42, 42}       % untuk blok peringatan
\definecolor{ubbgray}{RGB}{70, 76, 84}       % teks sekunder

% =====================================================================
%  IDENTITAS MATA KULIAH
% =====================================================================
\newcommand{\CourseName}{Pengantar Kriptografi}
\newcommand{\CourseCode}{MPFA014}
\newcommand{\CourseSemester}{Semester Ganjil}
\newcommand{\CourseProgram}{Program Studi Matematika}
\newcommand{\CourseFaculty}{Fakultas Matematika dan Ilmu Pengetahuan Alam}
\newcommand{\CourseInstitution}{Universitas Bale Bandung}
\newcommand{\CourseLecturer}{Cecep Suwanda, S.Si., M.Kom.}

% Label minggu, disetel tiap berkas lewat \setWeekInfo
\newcommand{\WeekLabel}{}
\newcommand{\WeekTopic}{}
\newcommand{\setWeekInfo}[2]{%
    \renewcommand{\WeekLabel}{#1}%
    \renewcommand{\WeekTopic}{#2}%
}

% =====================================================================
%  TEMA: header, footer, warna blok
% =====================================================================
\usetheme{default}
\setbeamertemplate{navigation symbols}{}
\setbeamertemplate{blocks}[rounded][shadow=false]

\setbeamercolor{palette primary}{bg=ubbblue, fg=white}
\setbeamercolor{palette secondary}{bg=ubbgray, fg=white}
\setbeamercolor{palette tertiary}{bg=ubbteal, fg=white}

\setbeamercolor{structure}{fg=ubbblue}
\setbeamercolor{normal text}{fg=black!92}
\setbeamercolor{alerted text}{fg=ubbred}

\setbeamercolor{block title}{bg=ubbblue, fg=white}
\setbeamercolor{block body}{bg=ubbblue!6, fg=black!92}
\setbeamercolor{block title alerted}{bg=ubbred, fg=white}
\setbeamercolor{block body alerted}{bg=ubbred!6, fg=black!92}
\setbeamercolor{block title example}{bg=ubbgreen, fg=white}
\setbeamercolor{block body example}{bg=ubbgreen!6, fg=black!92}

\setbeamercolor{itemize item}{fg=ubbblue}
\setbeamercolor{itemize subitem}{fg=ubbteal}
\setbeamercolor{enumerate item}{fg=ubbblue}

\setbeamerfont{headline}{size=\scriptsize}
\setbeamerfont{footline}{size=\scriptsize}
\setbeamerfont{block title}{size=\normalsize}

% Badan blok rata kiri (bukan rata kanan-kiri): teks slide lebih mudah dibaca
% dan baris terakhir yang pendek tidak lagi memicu "Underfull \hbox".
\addtobeamertemplate{block body begin}{\raggedright}{}
\addtobeamertemplate{block alerted body begin}{\raggedright}{}
\addtobeamertemplate{block example body begin}{\raggedright}{}

% Header: nama mata kuliah dan kode MK
\setbeamertemplate{headline}{%
    \begin{beamercolorbox}[wd=\paperwidth,ht=2.6ex,dp=1.2ex,leftskip=1em,rightskip=1em]{palette primary}%
        \usebeamerfont{headline}\CourseName\ \textperiodcentered\ \CourseCode%
        \hfill \CourseProgram%
    \end{beamercolorbox}%
}

% Footer: minggu dan nomor halaman
\setbeamertemplate{footline}{%
    \begin{beamercolorbox}[wd=\paperwidth,ht=2.4ex,dp=1.2ex,leftskip=1em,rightskip=1em]{palette secondary}%
        \usebeamerfont{footline}\WeekLabel\ \textperiodcentered\ \WeekTopic%
        \hfill \insertframenumber\,/\,\inserttotalframenumber%
    \end{beamercolorbox}%
}

% Halaman judul: judul, subjudul, lalu blok identitas mata kuliah
\setbeamertemplate{title page}{%
    \begin{centering}
        \vskip1.0em
        {\usebeamerfont{title}\usebeamercolor[fg]{title}\inserttitle\par}
        \vskip0.3em
        {\usebeamerfont{subtitle}\usebeamercolor[fg]{subtitle}\insertsubtitle\par}
        \vskip1.2em
        \textcolor{ubbblue}{\rule{0.72\textwidth}{0.8pt}}
        \vskip1.2em
        {\small
            \textbf{\CourseName}\ (\CourseCode)\par
            \vskip0.4em
            \CourseProgram\par
            \CourseFaculty\par
            \CourseInstitution\par
            \vskip0.6em
            \CourseSemester\ \textperiodcentered\ \WeekLabel\par
            \vskip0.8em
            Dosen Pengampu: \CourseLecturer\par
        }
        \vskip1.0em
        \textcolor{ubbblue}{\rule{0.72\textwidth}{0.4pt}}
    \end{centering}
}

% =====================================================================
%  PERINTAH BANTU
% =====================================================================
\newcommand{\R}{\mathbb{R}}
\newcommand{\N}{\mathbb{N}}
\newcommand{\Z}{\mathbb{Z}}
\newcommand{\code}[1]{\texttt{#1}}

% Rujukan gambar dari buku: \bookfig[lebar]{NN}{namaberkas.png}
% Jalur relatif terhadap folder ppt/ (dikompilasi dari sana).
\newcommand{\bookfig}[3][0.80\linewidth]{%
    \includegraphics[width=#1]{../buku_ajar/chapters/#2/figures/#3}%
}

% Keterangan sumber gambar, sepadan dengan \sumber di buku
\newcommand{\figsumber}[1]{%
    \par\vspace{1pt}{\tiny\itshape Sumber: #1\par}%
}

% Penanda blok materi yang sering dipakai di slide
\newcommand{\slidejudul}[1]{{\large\bfseries\textcolor{ubbblue}{#1}}\par\vspace{2pt}}
\newcommand{\kunci}[1]{\textcolor{ubbred}{\textbf{#1}}}

%      \title{...} \subtitle{...} \setWeekInfo{Minggu n}{Topik}
%
%  Berkas minggu hanya berisi \setWeekInfo dan lingkungan frame.
%  Kompilasi dijalankan DARI DALAM folder ppt/ (lihat compile-ppt.bat),
%  karena gambar buku dirujuk dengan jalur relatif ../buku_ajar/.
% =====================================================================

% ---------------------------------------------------------------- bahasa
\usepackage[utf8]{inputenc}
\usepackage[T1]{fontenc}
\usepackage[indonesian]{babel}
\usepackage{lmodern}
\usepackage{microtype}

% ------------------------------------------------------------- matematika
\usepackage{amsmath,amssymb,mathtools}

% ----------------------------------------------------------------- tabel
\usepackage{booktabs}
\usepackage{array}
\usepackage{tabularx}

% --------------------------------------------------------------- grafik
\usepackage{tikz}
\usetikzlibrary{arrows.meta,positioning,shapes.geometric,calc,fit,backgrounds}

% ---------------------------------------------------------- kode program
\usepackage{listings}
\lstdefinestyle{slidecode}{
    basicstyle=\ttfamily\scriptsize,
    keywordstyle=\color{ubbblue}\bfseries,
    commentstyle=\color{ubbgreen}\itshape,
    stringstyle=\color{ubbred},
    numbers=left,
    numberstyle=\tiny\color{gray},
    numbersep=6pt,
    frame=single,
    rulecolor=\color{ubbblue!40},
    backgroundcolor=\color{ubbblue!3},
    breaklines=true,
    showstringspaces=false,
    tabsize=4,
    columns=fullflexible,
    keepspaces=true,
}
\lstset{style=slidecode}

% =====================================================================
%  WARNA INSTITUSI
% =====================================================================
\definecolor{ubbblue}{RGB}{12, 62, 120}      % warna utama
\definecolor{ubbteal}{RGB}{0, 118, 110}      % aksen kedua
\definecolor{ubbgreen}{RGB}{22, 122, 74}     % untuk blok contoh
\definecolor{ubbred}{RGB}{168, 42, 42}       % untuk blok peringatan
\definecolor{ubbgray}{RGB}{70, 76, 84}       % teks sekunder

% =====================================================================
%  IDENTITAS MATA KULIAH
% =====================================================================
\newcommand{\CourseName}{Pengantar Kriptografi}
\newcommand{\CourseCode}{MPFA014}
\newcommand{\CourseSemester}{Semester Ganjil}
\newcommand{\CourseProgram}{Program Studi Matematika}
\newcommand{\CourseFaculty}{Fakultas Matematika dan Ilmu Pengetahuan Alam}
\newcommand{\CourseInstitution}{Universitas Bale Bandung}
\newcommand{\CourseLecturer}{Cecep Suwanda, S.Si., M.Kom.}

% Label minggu, disetel tiap berkas lewat \setWeekInfo
\newcommand{\WeekLabel}{}
\newcommand{\WeekTopic}{}
\newcommand{\setWeekInfo}[2]{%
    \renewcommand{\WeekLabel}{#1}%
    \renewcommand{\WeekTopic}{#2}%
}

% =====================================================================
%  TEMA: header, footer, warna blok
% =====================================================================
\usetheme{default}
\setbeamertemplate{navigation symbols}{}
\setbeamertemplate{blocks}[rounded][shadow=false]

\setbeamercolor{palette primary}{bg=ubbblue, fg=white}
\setbeamercolor{palette secondary}{bg=ubbgray, fg=white}
\setbeamercolor{palette tertiary}{bg=ubbteal, fg=white}

\setbeamercolor{structure}{fg=ubbblue}
\setbeamercolor{normal text}{fg=black!92}
\setbeamercolor{alerted text}{fg=ubbred}

\setbeamercolor{block title}{bg=ubbblue, fg=white}
\setbeamercolor{block body}{bg=ubbblue!6, fg=black!92}
\setbeamercolor{block title alerted}{bg=ubbred, fg=white}
\setbeamercolor{block body alerted}{bg=ubbred!6, fg=black!92}
\setbeamercolor{block title example}{bg=ubbgreen, fg=white}
\setbeamercolor{block body example}{bg=ubbgreen!6, fg=black!92}

\setbeamercolor{itemize item}{fg=ubbblue}
\setbeamercolor{itemize subitem}{fg=ubbteal}
\setbeamercolor{enumerate item}{fg=ubbblue}

\setbeamerfont{headline}{size=\scriptsize}
\setbeamerfont{footline}{size=\scriptsize}
\setbeamerfont{block title}{size=\normalsize}

% Badan blok rata kiri (bukan rata kanan-kiri): teks slide lebih mudah dibaca
% dan baris terakhir yang pendek tidak lagi memicu "Underfull \hbox".
\addtobeamertemplate{block body begin}{\raggedright}{}
\addtobeamertemplate{block alerted body begin}{\raggedright}{}
\addtobeamertemplate{block example body begin}{\raggedright}{}

% Header: nama mata kuliah dan kode MK
\setbeamertemplate{headline}{%
    \begin{beamercolorbox}[wd=\paperwidth,ht=2.6ex,dp=1.2ex,leftskip=1em,rightskip=1em]{palette primary}%
        \usebeamerfont{headline}\CourseName\ \textperiodcentered\ \CourseCode%
        \hfill \CourseProgram%
    \end{beamercolorbox}%
}

% Footer: minggu dan nomor halaman
\setbeamertemplate{footline}{%
    \begin{beamercolorbox}[wd=\paperwidth,ht=2.4ex,dp=1.2ex,leftskip=1em,rightskip=1em]{palette secondary}%
        \usebeamerfont{footline}\WeekLabel\ \textperiodcentered\ \WeekTopic%
        \hfill \insertframenumber\,/\,\inserttotalframenumber%
    \end{beamercolorbox}%
}

% Halaman judul: judul, subjudul, lalu blok identitas mata kuliah
\setbeamertemplate{title page}{%
    \begin{centering}
        \vskip1.0em
        {\usebeamerfont{title}\usebeamercolor[fg]{title}\inserttitle\par}
        \vskip0.3em
        {\usebeamerfont{subtitle}\usebeamercolor[fg]{subtitle}\insertsubtitle\par}
        \vskip1.2em
        \textcolor{ubbblue}{\rule{0.72\textwidth}{0.8pt}}
        \vskip1.2em
        {\small
            \textbf{\CourseName}\ (\CourseCode)\par
            \vskip0.4em
            \CourseProgram\par
            \CourseFaculty\par
            \CourseInstitution\par
            \vskip0.6em
            \CourseSemester\ \textperiodcentered\ \WeekLabel\par
            \vskip0.8em
            Dosen Pengampu: \CourseLecturer\par
        }
        \vskip1.0em
        \textcolor{ubbblue}{\rule{0.72\textwidth}{0.4pt}}
    \end{centering}
}

% =====================================================================
%  PERINTAH BANTU
% =====================================================================
\newcommand{\R}{\mathbb{R}}
\newcommand{\N}{\mathbb{N}}
\newcommand{\Z}{\mathbb{Z}}
\newcommand{\code}[1]{\texttt{#1}}

% Rujukan gambar dari buku: \bookfig[lebar]{NN}{namaberkas.png}
% Jalur relatif terhadap folder ppt/ (dikompilasi dari sana).
\newcommand{\bookfig}[3][0.80\linewidth]{%
    \includegraphics[width=#1]{../buku_ajar/chapters/#2/figures/#3}%
}

% Keterangan sumber gambar, sepadan dengan \sumber di buku
\newcommand{\figsumber}[1]{%
    \par\vspace{1pt}{\tiny\itshape Sumber: #1\par}%
}

% Penanda blok materi yang sering dipakai di slide
\newcommand{\slidejudul}[1]{{\large\bfseries\textcolor{ubbblue}{#1}}\par\vspace{2pt}}
\newcommand{\kunci}[1]{\textcolor{ubbred}{\textbf{#1}}}


\title{Kesepakatan Kunci di Saluran Terbuka}
\subtitle{Logaritma diskret, protokol Diffie--Hellman, dan algoritma ElGamal}
\setWeekInfo{Minggu 12}{Diffie--Hellman dan ElGamal}

\begin{document}

% ----------------------------------------------------------- halaman judul
\begin{frame}[plain]
    \titlepage
\end{frame}

% ------------------------------------------------------------------ peta
\begin{frame}{Peta Materi Pertemuan Ini}
    \begin{columns}[T,onlytextwidth]
        \begin{column}{0.52\textwidth}
            \slidejudul{Lima bagian}
            \footnotesize
            \begin{enumerate}
                \item \textbf{Logaritma diskret} --- akar primitif dan
                      persoalan membalik perpangkatan modular.
                \item \textbf{Diffie--Hellman} --- parameter publik, prosedur
                      empat langkah, dan bukti kesamaan kunci.
                \item \textbf{Serangan pada DH} --- penyadapan pasif dan
                      \textit{man-in-the-middle}.
                \item \textbf{ElGamal} --- pembangkitan kunci, enkripsi, dan
                      dekripsi.
                \item \textbf{Contoh terhitung} --- satu contoh DH dan satu
                      contoh ElGamal dari buku ajar.
            \end{enumerate}
        \end{column}
        \begin{column}{0.44\textwidth}
            \begin{block}{Referensi}
                \footnotesize
                Bab 11 dan Bab 12 buku ajar; Munir, \textit{IF4020
                Kriptografi}, ITB, 2026 (modul 15 dan 14)
            \end{block}
            \begin{exampleblock}{Benang merah}
                \footnotesize
                Diffie--Hellman membangun \kunci{kesepakatan kunci};
                ElGamal mengubah gagasan itu menjadi \kunci{enkripsi}.
                Keduanya bertumpu pada satu persoalan: logaritma diskret.
            \end{exampleblock}
        \end{column}
    \end{columns}
\end{frame}

% ---------------------------------------------------- capaian pembelajaran
\begin{frame}{Capaian Pembelajaran}
    \begin{block}{Sub-CPMK yang dibebankan}
        \textbf{Sub-CPMK 1.4} --- Mengenal algoritma kriptografi kunci publik,
        yaitu RSA, Diffie--Hellman, ElGamal, knapsack, dan \textit{Elliptic
        Curve Cryptography} (ECC), beserta persoalan matematis yang
        mendasarinya.
    \end{block}
    \vspace{0.2cm}
    \begin{exampleblock}{Setelah pertemuan ini, mahasiswa mampu}
        \footnotesize
        \begin{itemize}
            \item menjelaskan persoalan logaritma diskret dan akar primitif
                  sebagai landasan keamanan;
            \item menjelaskan prosedur pertukaran kunci Diffie--Hellman dan
                  alasan kedua pihak memperoleh kunci yang sama;
            \item mengidentifikasi serangan pasif dan serangan
                  \textit{man-in-the-middle} pada Diffie--Hellman;
            \item menjelaskan pembangkitan kunci, enkripsi, dan dekripsi
                  algoritma ElGamal;
            \item membandingkan peran Diffie--Hellman dan ElGamal dalam
                  kriptografi kunci publik.
        \end{itemize}
    \end{exampleblock}
\end{frame}

% =====================================================================
\section{Logaritma Diskret}
% =====================================================================

\begin{frame}{Dari Perpangkatan ke Logaritma Diskret}
    \begin{columns}[T,onlytextwidth]
        \begin{column}{0.54\textwidth}
            \footnotesize
            Perpangkatan modular mudah dihitung, bahkan untuk eksponen besar.
            Persoalan yang sukar adalah kebalikannya:
            \begin{itemize}
                \item diberikan $g$, $p$, dan $b$, carilah $x$ sehingga
                      $g^{x} \equiv b \pmod p$;
                \item nilai $x$ disebut \kunci{logaritma diskret} dari $b$
                      untuk basis $g$ modulo $p$;
                \item daerah jawabannya terbatas, yaitu $0 \le x \le p-1$.
            \end{itemize}
            \[
                b \equiv g^{x} \pmod p, \qquad 0 \le x \le p-1 .
            \]
        \end{column}
        \begin{column}{0.42\textwidth}
            \centering
            \begin{tikzpicture}[
                kotak/.style={draw, rounded corners, align=center,
                              font=\scriptsize, inner sep=4pt,
                              minimum width=3.2cm, minimum height=0.7cm},
                panah/.style={-{Latex[length=1.8mm]}, thick},
                balik/.style={-{Latex[length=1.8mm]}, thick, dashed,
                              ubbred},
                node distance=0.42cm,
            ]
                \node[kotak, fill=ubbblue!8] (maju)
                    {$g$, $x$, $p$};
                \node[kotak, fill=ubbgreen!10, below=of maju] (hasil)
                    {$b = g^{x} \bmod p$};
                \node[kotak, fill=ubbred!6, below=0.75cm of hasil] (balik)
                    {$x = \log_{g} b \pmod p$};
                \draw[panah] (maju) -- node[right, font=\tiny] {mudah} (hasil);
                \draw[balik] (hasil) -- node[right, font=\tiny] {sukar} (balik);
            \end{tikzpicture}
        \end{column}
    \end{columns}
    \vspace{0.1cm}
    \begin{alertblock}{Asimetri yang dimanfaatkan}
        \footnotesize
        Arah maju dapat diselesaikan dalam waktu polinomial terhadap
        $\log p$; arah balik tidak mempunyai cara efisien yang diketahui
        untuk $p$ yang besar.
    \end{alertblock}
\end{frame}

\begin{frame}{Akar Primitif: Syarat Basis yang Baik}
    \footnotesize
    Untuk bilangan prima $p$, bilangan $g$ disebut \kunci{akar primitif}
    bila perpangkatan $g^{1}, g^{2}, \dots, g^{p-1}$ modulo $p$
    menghasilkan seluruh residu $1$ sampai $p-1$ tepat satu kali.
    \vspace{0.15cm}
    \begin{columns}[T,onlytextwidth]
        \begin{column}{0.46\textwidth}
            \begin{block}{Contoh $p = 7$}
                \footnotesize
                $g = 3$ menghasilkan $(3,2,6,4,5,1)$, sehingga $3$ adalah
                akar primitif. $g = 2$ hanya menghasilkan $(2,4,1)$ yang
                berulang dengan periode $3$, sehingga bukan akar primitif.
            \end{block}
        \end{column}
        \begin{column}{0.50\textwidth}
            \centering
            \begin{tabular}{@{}l c c c c c c@{}}
                \toprule
                $i$ & 1 & 2 & 3 & 4 & 5 & 6 \\
                \midrule
                $3^{i} \bmod 7$ & 3 & 2 & 6 & 4 & 5 & 1 \\
                $2^{i} \bmod 7$ & 2 & 4 & 1 & 2 & 4 & 1 \\
                \bottomrule
            \end{tabular}
            \par\vspace{0.12cm}
            {\scriptsize Baris $g=3$ memuat semua residu bukan-nol.}
        \end{column}
    \end{columns}
    \vspace{0.1cm}
    \begin{alertblock}{Mengapa penting}
        \footnotesize
        Bila $g$ bukan akar primitif, tidak semua nilai $b$ mempunyai
        logaritma diskret. Basis yang salah membuat sebagian pesan tidak
        dapat diwakili.
    \end{alertblock}
\end{frame}

\begin{frame}{Contoh Logaritma Diskret}
    \begin{columns}[T,onlytextwidth]
        \begin{column}{0.54\textwidth}
            \footnotesize
            Bilangan $7$ adalah akar primitif dari $p = 41$. Kita mencari
            $x$ yang memenuhi
            \[
                15 \equiv 7^{x} \pmod{41} .
            \]
            Hitung pangkat secara bertahap:
            \[
                7^{1} = 7, \qquad
                7^{2} = 49 \equiv 8 \pmod{41},
            \]
            \[
                7^{3} = 343 = 8 \cdot 41 + 15 \equiv 15 \pmod{41} .
            \]
            Jadi, $x = 3$ adalah logaritma diskret dari $15$ untuk basis $7$
            modulo $41$.
        \end{column}
        \begin{column}{0.42\textwidth}
            \begin{exampleblock}{Cara memeriksa}
                \footnotesize
                Substitusikan kembali jawaban:
                \[
                    7^{3} = 343,
                \]
                \[
                    343 \bmod 41 = 15 .
                \]
                Hasilnya sama dengan nilai $b$ yang diberikan.
            \end{exampleblock}
            \begin{alertblock}{Skala persoalan}
                \footnotesize
                Untuk $p$ kecil, jawaban dapat dicari dengan mencoba satu per
                satu. Untuk $p$ besar, biaya pencarian menyeluruh tumbuh
                sebanding $p$.
            \end{alertblock}
        \end{column}
    \end{columns}
\end{frame}

\begin{frame}{Mengapa Logaritma Diskret Sukar}
    \footnotesize
    \begin{itemize}
        \item Ruang jawaban berukuran $p$: penelusuran menyeluruh menuntut
              sekitar $p$ percobaan.
        \item Perpangkatan modular dapat dihitung cepat dengan pengkuadratan
              berulang, jadi kecepatan arah maju tidak membantu arah balik.
        \item Tidak ada algoritma efisien yang diketahui untuk menghitung
              logaritma diskret pada $p$ besar secara umum.
        \item Karena itu $p$ harus dipilih sangat besar; nilai $p = 353$
              atau $p = 2357$ hanya dipakai untuk memperlihatkan langkah.
    \end{itemize}
    \vspace{0.1cm}
    \begin{exampleblock}{Akibat bagi dua algoritma}
        \footnotesize
        Diffie--Hellman dan ElGamal sama-sama menaruh rahasianya pada
        eksponen. Selama logaritma diskret sukar, eksponen itu tetap
        terlindungi meskipun nilai pangkatnya disiarkan.
    \end{exampleblock}
\end{frame}

% =====================================================================
\section{Pertukaran Kunci Diffie--Hellman}
% =====================================================================

\begin{frame}{Persoalan Distribusi Kunci}
    \begin{columns}[T,onlytextwidth]
        \begin{column}{0.56\textwidth}
            \footnotesize
            Cipher kunci-simetri menuntut kedua pihak sudah memegang kunci
            yang sama sebelum pesan pertama dikirim. Bila kunci itu harus
            dikirim lebih dahulu, saluran justru menjadi titik lemah.
            \begin{itemize}
                \item \textbf{Kriptografi hibrida}: kunci sesi dibungkus
                      dengan kunci publik penerima, lalu dipulihkan dengan
                      kunci privatnya.
                \item \textbf{Pertukaran kunci}: kedua pihak menghitung kunci
                      yang sama dari nilai-nilai yang dipertukarkan terbuka.
            \end{itemize}
        \end{column}
        \begin{column}{0.40\textwidth}
            \begin{block}{Diffie--Hellman (1976)}
                \footnotesize
                Kunci rahasia \emph{tidak pernah dikirim}. Yang melewati
                saluran hanyalah nilai publik yang tidak berguna bagi
                penyadap.
            \end{block}
            \begin{alertblock}{Batas tugasnya}
                \footnotesize
                DH tidak mengenkripsi pesan. Ia hanya menetapkan kunci
                bersama untuk dipakai cipher simetri.
            \end{alertblock}
        \end{column}
    \end{columns}
    \vspace{0.1cm}
    \begin{exampleblock}{Pertanyaan yang dijawab}
        \footnotesize
        Bagaimana Alice dan Bob sampai pada kunci yang sama tanpa seorang
        pun di antara mereka dapat menyimpulkan kunci itu dari lalu lintas
        yang terlihat?
    \end{exampleblock}
\end{frame}

\begin{frame}{Analogi Pencampuran Cat}
    \begin{columns}[T,onlytextwidth]
        \begin{column}{0.50\textwidth}
            \centering
            \bookfig[0.86\linewidth]{11}{analogi-cat-dh.png}
            \figsumber{Munir, \textit{IF4020 Kriptografi}, ITB.}
        \end{column}
        \begin{column}{0.48\textwidth}
            \footnotesize
            Alice dan Bob menyepakati warna publik, misalnya kuning.
            Masing-masing menambahkan warna rahasianya, bertukar campuran,
            lalu menambahkan lagi warna rahasianya sendiri.
            \begin{itemize}
                \item Pencampuran mudah dilakukan, tetapi memisahkan kembali
                      warna penyusunnya sukar.
                \item Urutan pencampuran tidak mengubah hasil akhir.
                \item Warna publik berperan seperti $p$ dan $g$; warna
                      rahasia berperan seperti eksponen $a$ dan $b$.
            \end{itemize}
            \begin{exampleblock}{Pesan analogi}
                \footnotesize
                Nilai yang dipertukarkan boleh dilihat semua orang, asalkan
                operasi baliknya cukup sukar.
            \end{exampleblock}
        \end{column}
    \end{columns}
\end{frame}

\begin{frame}{Parameter Publik dan Kunci Masing-Masing}
    \footnotesize
    Sebelum bertukar nilai, Alice dan Bob menyepakati dua parameter:
    \begin{itemize}
        \item $p$ adalah bilangan prima;
        \item $g$ adalah akar primitif dari $p$, dengan $g < p$.
    \end{itemize}
    Kedua nilai itu \kunci{tidak rahasia}; keduanya boleh dibicarakan
    melalui saluran terbuka.
    \vspace{0.15cm}
    \begin{columns}[T,onlytextwidth]
        \begin{column}{0.48\textwidth}
            \begin{block}{Alice}
                \footnotesize
                Memilih eksponen rahasia $a$, lalu menghitung kunci
                publiknya
                \[ A = g^{a} \bmod p . \]
            \end{block}
        \end{column}
        \begin{column}{0.48\textwidth}
            \begin{block}{Bob}
                \footnotesize
                Memilih eksponen rahasia $b$, lalu menghitung kunci
                publiknya
                \[ B = g^{b} \bmod p . \]
            \end{block}
        \end{column}
    \end{columns}
    \vspace{0.1cm}
    \begin{alertblock}{Yang rahasia dan yang publik}
        \footnotesize
        Hanya $a$ dan $b$ yang dirahasiakan. Nilai $p$, $g$, $A$, dan $B$
        boleh diketahui siapa saja tanpa membocorkan kunci bersama.
    \end{alertblock}
\end{frame}

\begin{frame}{Prosedur Pertukaran Kunci}
    \centering
    \begin{tikzpicture}[
        pihak/.style={draw, rounded corners, minimum width=2.6cm,
                      minimum height=0.62cm, font=\small},
        panah/.style={-{Latex[length=2.2mm]}, thick},
        label/.style={font=\scriptsize, align=center, fill=white,
                      inner sep=2pt},
        sisi/.style={font=\scriptsize, align=center, fill=white,
                     inner sep=2pt},
    ]
        \node[pihak, fill=ubbblue!8] (a) at (0,0) {Alice};
        \node[pihak, fill=ubbgreen!12] (b) at (7.2,0) {Bob};
        \draw[dashed] (a.south) -- ++(0,-4.05);
        \draw[dashed] (b.south) -- ++(0,-4.05);

        \node[sisi] at (0,-0.95) {1. pilih $a$ rahasia};
        \node[sisi] at (7.2,-0.95) {2. pilih $b$ rahasia};

        \draw[panah] (0,-1.85) -- node[label] {$A = g^{a} \bmod p$} (7.2,-1.85);
        \draw[panah] (7.2,-2.65) -- node[label] {$B = g^{b} \bmod p$} (0,-2.65);

        \node[sisi] at (0,-3.65) {3. $K = B^{a} \bmod p$};
        \node[sisi] at (7.2,-3.65) {4. $K = A^{b} \bmod p$};
    \end{tikzpicture}
    \begin{columns}[T,onlytextwidth]
        \begin{column}{0.58\textwidth}
            \footnotesize
            Alice mengirim $A$, Bob mengirim $B$. Masing-masing lalu
            memangkatkan nilai yang diterima dengan eksponen rahasianya
            sendiri.
        \end{column}
        \begin{column}{0.38\textwidth}
            \begin{exampleblock}{Hasil}
                \footnotesize
                Kedua pihak berakhir dengan satu nilai $K$ yang sama tanpa
                pernah mengirim $K$ itu.
            \end{exampleblock}
        \end{column}
    \end{columns}
\end{frame}

\begin{frame}{Mengapa Kedua Kunci Sama}
    \footnotesize
    Substitusikan definisi $A$ dan $B$ ke dalam perhitungan kedua pihak:
    \[
        \begin{aligned}
            K_{\text{Alice}}
                &= B^{a} \bmod p = (g^{b})^{a} \bmod p
                 = g^{ab} \bmod p, \\
            K_{\text{Bob}}
                &= A^{b} \bmod p = (g^{a})^{b} \bmod p
                 = g^{ab} \bmod p .
        \end{aligned}
    \]
    \vspace{0.1cm}
    \begin{columns}[T,onlytextwidth]
        \begin{column}{0.48\textwidth}
            \begin{block}{Sifat yang dipakai}
                \footnotesize
                $(g^{b})^{a} = g^{ba} = g^{ab} = (g^{a})^{b}$. Perkalian
                eksponen bersifat komutatif.
            \end{block}
        \end{column}
        \begin{column}{0.48\textwidth}
            \begin{exampleblock}{Urutan berbeda}
                \footnotesize
                Alice memangkatkan $B$ dengan $a$; Bob memangkatkan $A$
                dengan $b$. Kedua jalur bertemu pada residu yang sama.
            \end{exampleblock}
        \end{column}
    \end{columns}
\end{frame}

\begin{frame}{Sejarah dan Penamaan}
    \begin{columns}[T,onlytextwidth]
        \begin{column}{0.56\textwidth}
            \footnotesize
            \begin{itemize}
                \item Protokol ini dipublikasikan Whitfield Diffie dan
                      Martin Hellman pada tahun 1976.
                \item Istilah yang tepat adalah \emph{protokol}, bukan
                      algoritma, karena melibatkan pertukaran pesan antara
                      dua pihak.
                \item Gagasan distribusi kunci melalui saluran terbuka juga
                      dikembangkan Ralph Merkle.
                \item Pada tahun 2002 Hellman mengusulkan nama
                      \textit{Diffie--Hellman--Merkle} untuk mengakui
                      sumbangan Merkle.
                \item Ketiganya menerima IEEE Koji Kobayashi Computers and
                      Communications Award pada tahun 1999.
            \end{itemize}
        \end{column}
        \begin{column}{0.42\textwidth}
            \centering
            \bookfig[0.98\linewidth]{11}{pionir-diffie-hellman.png}
            \figsumber{Munir, \textit{IF4020 Kriptografi}, ITB.}
        \end{column}
    \end{columns}
\end{frame}

\begin{frame}{Ekstensi ke Tiga Pihak}
    \begin{columns}[T,onlytextwidth]
        \begin{column}{0.54\textwidth}
            \footnotesize
            Protokol dua pihak dapat diperluas kepada Alice, Bob, dan
            Carol yang menyepakati satu kunci bersama.
            \begin{itemize}
                \item Ketiganya menyepakati $p$ dan $g$, lalu memilih
                      eksponen rahasia $a$, $b$, dan $c$.
                \item Nilai diputar dalam lingkaran: setiap pihak
                      memangkatkan nilai yang diterima dengan eksponennya
                      sendiri, lalu meneruskannya.
                \item Setelah tiga jalur selesai, ketiganya memegang
                      $g^{abc} \bmod p$.
            \end{itemize}
        \end{column}
        \begin{column}{0.44\textwidth}
            \begin{block}{Tiga jalur menuju kunci yang sama}
                \footnotesize
                \[
                    (g^{ab})^{c} = (g^{bc})^{a} = (g^{ca})^{b}
                    = g^{abc} \pmod p .
                \]
                Urutan perkalian eksponen tidak mengubah hasil.
            \end{block}
            \begin{alertblock}{Bukan primitif baru}
                \footnotesize
                Ekstensi ini hanya mengulang gagasan DH dengan lebih banyak
                putaran dan pihak.
            \end{alertblock}
        \end{column}
    \end{columns}
\end{frame}

% =====================================================================
\section{Serangan pada Diffie--Hellman}
% =====================================================================

\begin{frame}{Penyadap Pasif: Berhadapan dengan Logaritma Diskret}
    \begin{columns}[T,onlytextwidth]
        \begin{column}{0.54\textwidth}
            \footnotesize
            Penyadap Eve melihat seluruh lalu lintas, yaitu
            \[
                p,\ g,\ A = g^{a} \bmod p,\ B = g^{b} \bmod p .
            \]
            Tidak satu pun nilai itu rahasia. Untuk memperoleh kunci, Eve
            harus lebih dahulu menemukan $a$ atau $b$:
            \[
                g^{a} \equiv A \pmod p, \qquad
                g^{b} \equiv B \pmod p .
            \]
            Keduanya adalah persoalan logaritma diskret.
        \end{column}
        \begin{column}{0.44\textwidth}
            \begin{block}{Bila eksponen diketahui}
                \footnotesize
                Dengan $a$, Eve menghitung $K = B^{a} \bmod p$. Dengan $b$,
                ia menghitung $K = A^{b} \bmod p$. Keduanya memberi kunci
                yang sama.
            \end{block}
            \begin{exampleblock}{Kesimpulan}
                \footnotesize
                Selama $p$ cukup besar, penyadapan pasif tidak memberi jalan
                menghitung $K$.
            \end{exampleblock}
        \end{column}
    \end{columns}
    \vspace{0.1cm}
    \begin{alertblock}{Batas keamanan}
        \footnotesize
        Keamanan DH terhadap penyadapan sepenuhnya bergantung pada
        kesukaran logaritma diskret. Bila $p$ kecil, seluruh rahasia dapat
        dihitung dengan penelusuran biasa.
    \end{alertblock}
\end{frame}

\begin{frame}{Serangan \textit{Man-in-the-Middle}: Mekanisme}
    \centering
    \begin{tikzpicture}[
        kotak/.style={draw, rounded corners, minimum height=0.8cm,
                      text width=2.1cm, align=center, font=\small,
                      inner sep=3pt},
        panah/.style={-{Latex[length=2.0mm]}, thick},
        palsu/.style={-{Latex[length=2.0mm]}, thick, dashed, ubbred},
        label/.style={font=\scriptsize, fill=white, inner sep=1pt},
    ]
        \node[kotak, fill=ubbblue!8]   (a) at (0,0)    {Alice};
        \node[kotak, fill=gray!10]     (m) at (5.5,0)  {Mallory};
        \node[kotak, fill=ubbgreen!12] (b) at (11,0)   {Bob};

        % arah Alice -- Mallory
        \draw[panah] (1.15,0.22)
            -- node[label, above] {$g^{a} \bmod p$} (4.35,0.22);
        \draw[palsu] (4.35,-0.22)
            -- node[label, below] {$g^{m} \bmod p$} (1.15,-0.22);

        % arah Mallory -- Bob
        \draw[palsu] (6.65,0.22)
            -- node[label, above] {$g^{m} \bmod p$} (9.85,0.22);
        \draw[panah] (9.85,-0.22)
            -- node[label, below] {$g^{b} \bmod p$} (6.65,-0.22);

        \node[font=\scriptsize, align=center] at (0,-1.2)
            {memegang $a$};
        \node[font=\scriptsize, align=center] at (5.5,-1.2)
            {memegang $m$};
        \node[font=\scriptsize, align=center] at (11,-1.2)
            {memegang $b$};
    \end{tikzpicture}
    \vspace{0.15cm}
    \begin{columns}[T,onlytextwidth]
        \begin{column}{0.56\textwidth}
            \footnotesize
            Mallory menyamar sebagai Bob di hadapan Alice, dan sebagai
            Alice di hadapan Bob. Ia mencegat $g^{a}$ dan $g^{b}$, lalu
            menggantikan keduanya dengan $g^{m}$.
        \end{column}
        \begin{column}{0.40\textwidth}
            \begin{alertblock}{Yang tidak perlu diketahui}
                \footnotesize
                Mallory tidak perlu menemukan $a$ atau $b$. Ia hanya
                memakai eksponen rahasianya sendiri, $m$.
            \end{alertblock}
        \end{column}
    \end{columns}
\end{frame}

\begin{frame}{Serangan \textit{Man-in-the-Middle}: Akibat}
    \footnotesize
    Setelah pertukaran yang dipalsukan, kedua pihak menghitung kunci yang
    berbeda:
    \[
        K_1 = (g^{m})^{a} \bmod p = g^{am} \bmod p,
        \qquad
        K_2 = (g^{m})^{b} \bmod p = g^{bm} \bmod p .
    \]
    Mallory dapat menghitung keduanya dari nilai publik:
    \[
        K_1 = A^{m} \bmod p, \qquad K_2 = B^{m} \bmod p .
    \]
    \begin{itemize}
        \item Alice mengira $K_1$ adalah kunci dengan Bob; Bob mengira
              $K_2$ adalah kunci dengan Alice.
        \item Mallory membuka pesan dengan $K_1$, lalu mengenkripsinya
              kembali dengan $K_2$ sebelum diteruskan.
        \item Alice dan Bob tidak menyadari bahwa mereka tidak pernah
              berbagi kunci yang sama.
    \end{itemize}
    \begin{alertblock}{Pelajaran}
        \footnotesize
        DH aman terhadap penyadapan pasif, tetapi tidak menahan serangan
        aktif. Kesepakatan kunci perlu diikat pada identitas pihak yang
        berkomunikasi.
    \end{alertblock}
\end{frame}

% =====================================================================
\section{Algoritma ElGamal}
% =====================================================================

\begin{frame}{ElGamal: dari DH ke Enkripsi Kunci Publik}
    \begin{columns}[T,onlytextwidth]
        \begin{column}{0.62\textwidth}
            \footnotesize
            \begin{itemize}
                \item Algoritma ElGamal diperkenalkan Taher Elgamal pada
                      tahun 1985.
                \item Makalahnya terbit di \textit{IEEE Transactions on
                      Information Theory}, volume 31, nomor 4.
                \item Keamanannya bersandar pada kesukaran menghitung
                      logaritma diskret di lapangan berhingga.
                \item Gagasan utamanya: skema distribusi kunci
                      Diffie--Hellman diubah menjadi sistem kriptografi
                      kunci publik.
                \item RSA bertumpu pada kesukaran pemfaktoran; ElGamal
                      bertumpu pada logaritma diskret.
            \end{itemize}
        \end{column}
        \begin{column}{0.34\textwidth}
            \centering
            \bookfig[0.94\linewidth]{12}{taher-elgamal.png}
            \figsumber{Munir, \textit{IF4020 Kriptografi}, ITB.}
        \end{column}
    \end{columns}
    \vspace{0.1cm}
    \begin{exampleblock}{Perbedaan peran}
        \footnotesize
        DH menghasilkan satu kunci untuk dua pihak. ElGamal memakai
        kunci publik seorang penerima untuk mengenkripsi pesan, sehingga
        pengirim tidak perlu berbagi rahasia lebih dahulu.
    \end{exampleblock}
\end{frame}

\begin{frame}{Parameter ElGamal dan Status Kerahasiaan}
    \footnotesize
    \begin{tabularx}{\linewidth}{@{}l l X c@{}}
        \toprule
        \textbf{Besaran} & \textbf{Arti} & \textbf{Peran} & \textbf{Rahasia?} \\
        \midrule
        $p$ & bilangan prima & modulus publik & tidak \\
        \addlinespace
        $g$ & akar primitif dari $p$ & basis publik & tidak \\
        \addlinespace
        $x$ & eksponen, $2 \le x \le p-2$ & kunci privat & \textbf{ya} \\
        \addlinespace
        $y$ & $y = g^{x} \bmod p$ & kunci publik & tidak \\
        \addlinespace
        $m$ & plainteks, $0 \le m \le p-1$ & pesan & \textbf{ya} \\
        \addlinespace
        $a,b$ & pasangan cipherteks & dikirim & tidak \\
        \bottomrule
    \end{tabularx}
    \vspace{0.2cm}
    \begin{columns}[T,onlytextwidth]
        \begin{column}{0.48\textwidth}
            \begin{block}{Kunci}
                \footnotesize
                Kunci publik berupa $(y,g,p)$. Kunci privat berupa
                $(x,p)$; nilai $x$ tidak pernah meninggalkan pemiliknya.
            \end{block}
        \end{column}
        \begin{column}{0.48\textwidth}
            \begin{alertblock}{Sumber keamanan}
                \footnotesize
                Penyerang yang melihat $y$ harus menemukan $x$ dari
                $y = g^{x} \bmod p$. Itu persoalan logaritma diskret.
            \end{alertblock}
        \end{column}
    \end{columns}
\end{frame}

\begin{frame}{Pembangkitan Pasangan Kunci}
    \footnotesize
    Pemilik kunci menjalankan tiga langkah:
    \begin{enumerate}
        \item pilih bilangan prima $p$;
        \item pilih $g$ yang merupakan akar primitif dari $p$, lalu pilih
              kunci privat $x$ dengan $2 \le x \le p-2$;
        \item hitung kunci publik
              \[
                  y = g^{x} \bmod p .
              \]
    \end{enumerate}
    \vspace{0.1cm}
    \begin{columns}[T,onlytextwidth]
        \begin{column}{0.46\textwidth}
            \begin{exampleblock}{Yang disiarkan}
                \footnotesize
                Tripel $(y,g,p)$ boleh diumumkan. Siapa pun dapat memakainya
                untuk mengirim pesan kepada pemilik kunci.
            \end{exampleblock}
        \end{column}
        \begin{column}{0.50\textwidth}
            \begin{alertblock}{Yang disimpan}
                \footnotesize
                Pasangan $(x,p)$ tetap rahasia. Dari $y$ dan $g$, nilai $x$
                hanya dapat diperoleh dengan menyelesaikan logaritma
                diskret.
            \end{alertblock}
        \end{column}
    \end{columns}
\end{frame}

\begin{frame}{Enkripsi ElGamal}
    \footnotesize
    Pengirim memakai kunci publik $(y,g,p)$ dan menjalankan tiga langkah:
    \begin{enumerate}
        \item pecah plainteks menjadi blok $m_1, m_2, \dots$ dengan syarat
              setiap blok berada pada selang $0 \le m \le p-1$;
        \item pilih bilangan acak $k$ dengan $1 \le k \le p-2$;
        \item untuk setiap blok $m$, hitung pasangan
              \[
                  a = g^{k} \bmod p,
                  \qquad
                  b = y^{k} \cdot m \bmod p .
              \]
    \end{enumerate}
    \vspace{0.1cm}
    \begin{columns}[T,onlytextwidth]
        \begin{column}{0.48\textwidth}
            \begin{block}{Cipherteks}
                \footnotesize
                Pasangan $(a,b)$ adalah cipherteks satu blok. Ukurannya dua
                kali ukuran plainteks karena terdiri atas dua residu modulo
                $p$.
            \end{block}
        \end{column}
        \begin{column}{0.48\textwidth}
            \begin{alertblock}{Peran $k$}
                \footnotesize
                $k$ dipilih acak untuk setiap pesan. Nilai $a = g^{k}$
                tidak bergantung pada $m$; pesan disamarkan oleh faktor
                $y^{k}$.
            \end{alertblock}
        \end{column}
    \end{columns}
\end{frame}

\begin{frame}{Dekripsi ElGamal}
    \footnotesize
    Penerima memakai kunci privat $x$ untuk memulihkan $m$ dari $(a,b)$.
    \begin{enumerate}
        \item hitung invers dari $a^{x}$ modulo $p$:
              \[
                  (a^{x})^{-1} \equiv a^{\,p-1-x} \pmod p ;
              \]
        \item kalikan cipherteks $b$ dengan invers itu:
              \[
                  m = b \cdot (a^{x})^{-1} \bmod p .
              \]
    \end{enumerate}
    \vspace{0.1cm}
    \begin{columns}[T,onlytextwidth]
        \begin{column}{0.50\textwidth}
            \begin{block}{Mengapa rumus invers benar}
                \footnotesize
                $a^{x} \cdot a^{p-1-x} = a^{p-1} \equiv 1 \pmod p$,
                sehingga $a^{p-1-x}$ memang invers dari $a^{x}$.
            \end{block}
        \end{column}
        \begin{column}{0.46\textwidth}
            \begin{exampleblock}{Yang dihitung penerima}
                \footnotesize
                Hanya satu perpangkatan untuk invers dan satu perkalian
                modular untuk memulihkan pesan.
            \end{exampleblock}
        \end{column}
    \end{columns}
\end{frame}

\begin{frame}{Mengapa Dekripsi Berhasil}
    \footnotesize
    Substitusikan definisi $b$, $y$, dan $a$:
    \[
        \begin{aligned}
            \frac{b}{a^{x}}
                &\equiv \frac{y^{k} m}{a^{x}} \pmod p \\
                &\equiv \frac{g^{xk} m}{g^{xk}} \pmod p \\
                &\equiv m \pmod p .
        \end{aligned}
    \]
    \vspace{0.1cm}
    \begin{columns}[T,onlytextwidth]
        \begin{column}{0.48\textwidth}
            \begin{block}{Dua kesamaan yang dipakai}
                \footnotesize
                $a^{x} = (g^{k})^{x} = g^{kx}$ dan
                $y^{k} = (g^{x})^{k} = g^{xk}$.
            \end{block}
        \end{column}
        \begin{column}{0.48\textwidth}
            \begin{exampleblock}{Faktor yang saling meniadakan}
                \footnotesize
                Faktor $g^{xk}$ muncul di pembilang dan penyebut. Yang
                tersisa hanya $m$.
            \end{exampleblock}
        \end{column}
    \end{columns}
\end{frame}

\begin{frame}{Sifat Probabilistik dan Ukuran Cipherteks}
    \footnotesize
    \begin{itemize}
        \item Karena $k$ dipilih acak, pesan yang sama menghasilkan
              pasangan $(a,b)$ yang berbeda pada setiap enkripsi.
        \item $k$ harus dipilih baru untuk setiap pesan; nilai acak itu
              adalah bagian dari algoritma, bukan tambahan opsional.
        \item Cipherteks berukuran dua kali plainteks karena memuat dua
              residu, yaitu $a$ dan $b$.
        \item Komponen $a = g^{k}$ tidak memuat $m$; informasi pesan
              tersembunyi pada $b = y^{k}m \bmod p$.
        \item Dekripsi menghapus topeng $y^{k}$ dengan memakai $x$.
    \end{itemize}
    \vspace{0.15cm}
    \begin{exampleblock}{Akibat praktis}
        \footnotesize
        Dua pengiriman pesan yang identik tidak menghasilkan cipherteks
        yang identik. Sifat ini berbeda dari cipher deterministik yang
        selalu memberi keluaran sama untuk masukan yang sama.
    \end{exampleblock}
\end{frame}

% =====================================================================
\section{Contoh Terhitung}
% =====================================================================

\begin{frame}{Contoh Diffie--Hellman (1/2): Kunci Publik}
    \footnotesize
    Alice dan Bob menyepakati $p = 353$ dan $g = 3$. Bilangan $3$ adalah
    akar primitif dari $353$, dan syarat $g < p$ terpenuhi.
    \vspace{0.15cm}
    \begin{block}{Langkah 1 --- Alice memilih $a = 97$ dan menghitung}
        \footnotesize
        \[
            A = g^{a} \bmod p = 3^{97} \bmod 353 = 40 .
        \]
    \end{block}
    \begin{block}{Langkah 2 --- Bob memilih $b = 233$ dan menghitung}
        \footnotesize
        \[
            B = g^{b} \bmod p = 3^{233} \bmod 353 = 248 .
        \]
    \end{block}
    \begin{exampleblock}{Yang melewati saluran}
        \footnotesize
        Alice mengirim $A = 40$ kepada Bob; Bob mengirim $B = 248$ kepada
        Alice. Kedua nilai itu publik.
    \end{exampleblock}
\end{frame}

\begin{frame}{Contoh Diffie--Hellman (2/2): Kunci Bersama}
    \footnotesize
    Setelah bertukar kunci publik, kedua pihak menghitung kunci rahasia
    dari sisi masing-masing.
    \vspace{0.15cm}
    \begin{block}{Langkah 3 --- Alice memakai $B = 248$ dan $a = 97$}
        \footnotesize
        \[
            K = B^{a} \bmod p = 248^{97} \bmod 353 = 160 .
        \]
    \end{block}
    \begin{block}{Langkah 4 --- Bob memakai $A = 40$ dan $b = 233$}
        \footnotesize
        \[
            K = A^{b} \bmod p = 40^{233} \bmod 353 = 160 .
        \]
    \end{block}
    \begin{exampleblock}{Pemeriksaan dan kesimpulan}
        \footnotesize
        Kedua jalur bertemu pada $K = 160$, sama dengan
        $g^{ab} = 3^{97 \cdot 233} \bmod 353 = 160$. Kunci ini siap dipakai
        cipher simetri.
    \end{exampleblock}
\end{frame}

\begin{frame}{Contoh ElGamal (1/3): Pembangkitan Kunci}
    \footnotesize
    Bob memilih parameter dan kunci privatnya:
    \[
        p = 2357, \qquad g = 2, \qquad x = 1751 .
    \]
    Syarat terpenuhi: $2$ adalah akar primitif dari $2357$ dan
    $2 \le 1751 \le 2355$.
    \vspace{0.15cm}
    \begin{block}{Menghitung kunci publik}
        \footnotesize
        \[
            y = g^{x} \bmod p = 2^{1751} \bmod 2357 = 1185 .
        \]
    \end{block}
    \begin{columns}[T,onlytextwidth]
        \begin{column}{0.48\textwidth}
            \begin{exampleblock}{Kunci publik Bob}
                \footnotesize
                $(y,g,p) = (1185, 2, 2357)$. Nilai ini diberitahukan kepada
                Alice.
            \end{exampleblock}
        \end{column}
        \begin{column}{0.48\textwidth}
            \begin{alertblock}{Kunci privat Bob}
                \footnotesize
                $(x,p) = (1751, 2357)$. Nilai $x$ tidak dikirim kepada
                siapa pun.
            \end{alertblock}
        \end{column}
    \end{columns}
\end{frame}

\begin{frame}{Contoh ElGamal (2/3): Enkripsi}
    \footnotesize
    Alice akan mengirim pesan $m = 2035$. Nilai itu berada di dalam selang
    $[0,2356]$, sehingga dapat dienkripsi sebagai satu blok.
    \vspace{0.1cm}
    \begin{block}{Alice memilih $k = 1520$ dan menghitung}
        \footnotesize
        \[
            a = g^{k} \bmod p = 2^{1520} \bmod 2357 = 1430,
        \]
        \[
            b = y^{k} \cdot m \bmod p
              = 1185^{1520} \cdot 2035 \bmod 2357 = 697 .
        \]
    \end{block}
    \begin{exampleblock}{Cipherteks yang dikirim}
        \footnotesize
        Pasangan $(a,b) = (1430, 697)$. Ukurannya dua kali ukuran
        plainteks, sesuai sifat ElGamal.
    \end{exampleblock}
\end{frame}

\begin{frame}{Contoh ElGamal (3/3): Dekripsi}
    \footnotesize
    Bob memakai kunci privat $(x,p) = (1751, 2357)$ pada cipherteks
    $(a,b) = (1430,697)$.
    \vspace{0.1cm}
    \begin{block}{Langkah 1 --- menghitung invers}
        \footnotesize
        Eksponen invers adalah $p-1-x = 2357-1-1751 = 605$:
        \[
            (a^{x})^{-1} = a^{\,p-1-x} \bmod p
            = 1430^{605} \bmod 2357 = 872 .
        \]
    \end{block}
    \begin{block}{Langkah 2 --- memulihkan pesan}
        \footnotesize
        \[
            m = b \cdot (a^{x})^{-1} \bmod p
              = 697 \cdot 872 \bmod 2357 = 2035 .
        \]
    \end{block}
    \begin{exampleblock}{Pemeriksaan}
        \footnotesize
        $a^{x} = 1430^{1751} \bmod 2357 = 2084$, dan
        $2084 \cdot 872 \bmod 2357 = 1$. Jadi $872$ memang invers yang
        benar, dan plainteks $m = 2035$ kembali utuh.
    \end{exampleblock}
\end{frame}

\begin{frame}{Contoh ElGamal: Pesan ``HALO'' Multi-Blok}
    \footnotesize
    Dengan pengodean $A = 00, \dots, Z = 25$, kata ``HALO'' menjadi
    \[
        \mathrm{H}=07,\quad \mathrm{A}=00,\quad \mathrm{L}=11,\quad
        \mathrm{O}=14
        \qquad \Longrightarrow \qquad 07001114 .
    \]
    Parameter Alice: $p = 2273$, $g = 3$, $x = 243$, dan
    $y = 3^{243} \bmod 2273 = 461$.

    Pesan dipecah menjadi $m_1 = 0700 = 700$ dan $m_2 = 1114$.
    \vspace{0.1cm}
    \centering
    \begin{tabular}{@{}l c c c c@{}}
        \toprule
        Blok & $m$ & $k$ & $a = g^{k} \bmod p$ & $b = y^{k}m \bmod p$ \\
        \midrule
        $m_1$ & 700  & 1463 & $3^{1463} \bmod 2273 = 1439$
              & $461^{1463} \cdot 700 \bmod 2273 = 74$ \\
        \addlinespace
        $m_2$ & 1114 & 2001 & $3^{2001} \bmod 2273 = 1220$
              & $461^{2001} \cdot 1114 \bmod 2273 = 1682$ \\
        \bottomrule
    \end{tabular}
    \vspace{0.12cm}
    \begin{exampleblock}{Dekripsi dengan $x = 243$}
        \footnotesize
        $74 \cdot 1439^{2029} \bmod 2273 = 700$ dan
        $1682 \cdot 1220^{2029} \bmod 2273 = 1114$. Rangkaian
        $07001114$ kembali berbunyi ``HALO''.
    \end{exampleblock}
\end{frame}

% =====================================================================
\section{Rangkuman}
% =====================================================================

\begin{frame}{Perbandingan Diffie--Hellman dan ElGamal}
    \footnotesize
    \begin{tabularx}{\linewidth}{@{}>{\bfseries\raggedright\arraybackslash}p{2.3cm}
                                   >{\raggedright\arraybackslash}X
                                   >{\raggedright\arraybackslash}X@{}}
        \toprule
        Aspek & Diffie--Hellman & ElGamal \\
        \midrule
        Tujuan & menyepakati satu kunci bersama & mengenkripsi pesan dengan
        kunci publik \\
        \addlinespace
        Keluaran & satu nilai $K = g^{ab} \bmod p$ & pasangan cipherteks
        $(a,b)$ \\
        \addlinespace
        Rahasia & eksponen $a$ dan $b$ & kunci privat $x$ dan bilangan acak
        $k$ \\
        \addlinespace
        Publik & $p$, $g$, $A$, $B$ & $p$, $g$, $y$, dan $(a,b)$ \\
        \addlinespace
        Landasan & logaritma diskret & logaritma diskret \\
        \addlinespace
        Sifat & protokol dua arah & enkripsi probabilistik \\
        \bottomrule
    \end{tabularx}
    \vspace{0.2cm}
    \begin{exampleblock}{Hubungan keduanya}
        \footnotesize
        ElGamal dapat dipandang sebagai gagasan DH yang dipakai untuk
        mengenkripsi: kunci publik $y$ berperan seperti nilai publik DH,
        sedangkan kunci privat $x$ berperan seperti eksponen rahasianya.
    \end{exampleblock}
\end{frame}

\begin{frame}{Rangkuman}
    \footnotesize
    \begin{itemize}
        \item Logaritma diskret adalah kebalikan perpangkatan modular:
              diberikan $g$, $p$, dan $b$, carilah $x$ pada
              $g^{x} \equiv b \pmod p$.
        \item Arah maju mudah dihitung; arah balik sukar untuk $p$ besar.
              Inilah satu-satunya sumber keamanan DH dan ElGamal.
        \item DH memakai $p$ dan $g$ publik, eksponen $a$ dan $b$ rahasia,
              lalu bertukar $A = g^{a}$ dan $B = g^{b}$.
        \item Kedua pihak memperoleh $K = g^{ab} \bmod p$ karena perkalian
              eksponen komutatif.
        \item DH aman terhadap penyadapan pasif, tetapi tidak menahan
              \textit{man-in-the-middle} yang mengganti $A$ dan $B$ dengan
              $g^{m}$.
        \item ElGamal memakai kunci publik $(y,g,p)$ dan kunci privat
              $(x,p)$; enkripsi memberi $a = g^{k}$ dan
              $b = y^{k}m \bmod p$.
        \item Dekripsi berhasil karena $y^{k} = a^{xk}$, sehingga faktor
              itu saling meniadakan dan menyisakan $m$.
    \end{itemize}
\end{frame}

\begin{frame}{Penutup}
    \begin{block}{Yang harus diingat}
        \footnotesize
        \begin{itemize}
            \item Dasar keamanan kedua algoritma adalah eksponen rahasia;
                  nilai pangkatnya boleh disiarkan.
            \item Diffie--Hellman menetapkan kunci, bukan mengirimkannya.
            \item Tanpa autentikasi pihak, DH dapat dibelokkan oleh
                  \textit{man-in-the-middle}.
            \item ElGamal memperluas DH menjadi enkripsi dengan pasangan
                  $(a,b)$ dan bilangan acak $k$.
            \item Cipherteks ElGamal dua kali lebih panjang dan berubah
                  setiap kali pesan dienkripsi.
        \end{itemize}
    \end{block}
    \vspace{0.15cm}
    \begin{exampleblock}{Pertemuan berikutnya}
        \textbf{Minggu 13 --- Kriptografi kurva eliptik}: grup dan medan
        berhingga seperlunya, kurva eliptik, ECDLP, ECDH, ECEG, keamanan,
        dan keuntungan ECC.
    \end{exampleblock}
\end{frame}

\end{document}
